% บทคัดย่อภาษาไทย ใช้ใน main_short_th.tex (ฉบับอังกฤษอยู่ใน 00_abstract.tex); ตัวเลขทั้งหมดมาจาก numbers.tex.
% environment thaiabstract นิยามใน main_short_th.tex; ต้อง build ด้วย XeTeX/tectonic.
\begin{thaiabstract}
โทรศัพท์ส่วนใหญ่ไม่มี LiDAR จึงต้องสแกนห้องด้วยการประมาณความลึกจากภาพเดี่ยว
แต่โมเดลสำเร็จรูปที่ให้ความลึกเป็นเมตร เช่น Depth Anything V2 (DA-v2) Metric-Indoor~\cite{da2}
ซึ่งฝึกด้วยภาพสังเคราะห์ ทำให้ห้องที่สร้างได้คลาดเคลื่อน \rsMETRICCM{}~ซม.
งานวิจัยนี้ศึกษาว่าการปรับจูน (fine-tune) โมเดลดังกล่าวด้วยความลึกจาก LiDAR ของโทรศัพท์เอง
โดยไม่ใช้เครื่องสแกนเลเซอร์ ทำให้แม่นยำกว่าโมเดลเดิมหรือไม่ เมื่อวัดเทียบกับเครื่องสแกนเลเซอร์
Faro Focus S70~\cite{faro} ซึ่งเป็นค่าอ้างอิงของชุดข้อมูล ARKitScenes~\cite{arkitscenes}
โดยปรับจูนเฉพาะส่วน DPT head~\cite{dpt} (ตัวถอดรหัสความลึก โดยตรึงส่วนเข้ารหัสภาพไว้)
ด้วยความลึกจาก LiDAR ของ iPad บน 24 ห้อง ได้โมเดล \ftmain{}
แล้วเปรียบเทียบกับโมเดลเดิมในกระบวนการเดียวกัน คือรวมความลึกรายเฟรมเป็นปริมาตร
TSDF (truncated signed distance function)~\cite{curless96} แล้วสร้างเป็นพื้นผิวสามมิติ
โดยให้แหล่งความลึกเป็นตัวแปรเดียว บนห้องที่ไม่เคยใช้ฝึก ได้แก่ 6 ห้องที่มีแบบจำลองอ้างอิงจากเครื่องสแกนเลเซอร์
และอีก \rsVN{} ห้องที่ประเมินรายเฟรม
ผลคือ \ftmain{} ลดค่า Chamfer distance~\cite{chamfer}
(ระยะเฉลี่ยถึงจุดที่ใกล้ที่สุดระหว่างพื้นผิวที่สร้างได้กับพื้นผิวอ้างอิง วัดทั้งสองทิศ) จาก \rsMETRICCM{}~ซม. เหลือ \rsFTLIDARALLCM{}~ซม.
(ดีขึ้น \rsMAINGAIN{} เท่า และดีขึ้นทุกห้อง) และลดค่า AbsRel~\cite{eigen14}
(ความคลาดเคลื่อนสัมพัทธ์เฉลี่ยของความลึก) จาก \rsVMETRICABSREL{} เหลือ \rsVFTALLABSREL{}
(ดีขึ้น \rsVWINS{} จาก \rsVN{} ห้อง) โดยไม่ต้องปรับเทียบหรือใช้เซนเซอร์ขณะใช้งาน
และการใช้เครื่องสแกนเลเซอร์เป็นครูไม่ได้ดีกว่าอย่างวัดได้
ความคลาดเคลื่อนที่เหลือมาจากสเกลที่แกว่งตามมุมมองของแต่ละเฟรม ซึ่งการปรับจูนขจัดไม่ได้
จึงยังด้อยกว่า LiDAR ของ iPad (\rsLIDARCM{}~ซม.) และยังไม่ถึงระดับความแม่นยำสำหรับดูภาพรวม
(พื้นผิวของห้องอย่างน้อย 90\,\% อยู่ภายใน 10~ซม.)
\end{thaiabstract}
